What PS represents, its five weighted components plus the conditional CoT bonus, the per-row bootstrap CI, how to read the capability radar, and how to retune the weights for a specific research goal.

The Composite Planning Score collapses five orthogonal capability axes into one number per model (and per model\(\times\)domain). Its purpose is triage, not verdict: a single ranking to scan, always reported alongside its components because a composite can hide compensating strengths and weaknesses.

\subsubsection{Why a composite at all}

The preceding sections give nine-plus dials. PS exists so that, for a quick comparison, you have one axis that already encodes the framework's value judgement: zero-shot ability matters most, retries are a cost, structural quality and grounding matter, and precondition awareness is a tie-breaker. It is the metric you sort the leaderboard by; it is never the metric you conclude from.

\subsubsection{The components and their weights}

PS is a weighted sum of five inputs, each already in [0,1], with an optional sixth (CoT) bonus. The person sometimes thinks of it as "two components" --- that intuition maps onto the structure exactly: there is a \textbf{base}(the five-term weighted sum) and a \textbf{conditional CoT bonus} that only applies when a strict CoT alignment score exists, at which point all weights are renormalised so the result stays bounded.

\begin{bmformulabox}
\begin{align*}
\text{base} &=
0.25 \cdot \text{FASR}
+ 0.20 \cdot \text{IWSR}
+ 0.20 \cdot \text{exec\_ratio}
+ 0.20 \cdot (1 - \text{halluc})
+ 0.15 \cdot \text{PAS}
\\[4pt]
\text{PS} &=
\frac{\text{base} + 0.05 \cdot \text{CoT}}{1.05},
\quad \text{if CoT is finite}
\\[4pt]
&\rightarrow \text{renormalised and clipped to } [0,1]
\end{align*}
\end{bmformulabox}

\begin{table}[h]\small
\centering
\begin{tabularx}{\linewidth}{@{}l r I P@{}}
\toprule
Component & w & Rationale for the weight & NaN substitute\\
\midrule
FASR & 0.25 & Highest --- the only signal that cannot be inflated by retries & 0\\
\addlinespace
IWSR & 0.20 & Retry efficiency; rewards fast convergence, penalises stochastic search & 0\\
\addlinespace
exec\_ratio & 0.20 & Structural quality independent of goal success & 0\\
\addlinespace
1 - halluc & 0.20 & Schema grounding & 0.5\\
\addlinespace
PAS & 0.15 & Lowest --- undefined for zero-failure models, so a 0.5 prior is needed & 0.5\\
\addlinespace
CoT bonus & +0.05 & Applied only when strict CoT alignment is finite & skipped\\
\bottomrule
\end{tabularx}
\end{table}

\subsubsection{Aggregation \& range}

PS is computed on the \emph{aggregated} components in both the \code{overall} and \code{by\_domain} tables --- so each model carries an overall PS and one PS per domain. In the bundled run: gpt\_oss 0.548 \(\cdot\) nemotron(nv) 0.446 \(\cdot\) kimi 0.435 \(\cdot\) qwen\_3\_5 0.425, down to phi\_4\_mini 0.315.

\subsubsection{The 95\% bootstrap confidence interval}

To show whether a PS difference is real or noise, the \code{composite\_score} table attaches a bootstrap CI (N=1000, seed=42). Crucially it resamples \textbf{per-row PS contributions}, not per-model means --- a more conservative estimate that reflects instance-level variability.

\begin{lstlisting}
contribs = model_df.apply(per_row_composite, axis=1).values
boots = [contribs[rng.integers(0,len(contribs),len(contribs))].mean() for _ in range(1000)]
PS_ci_lo, PS_ci_hi = np.percentile(boots,2.5), np.percentile(boots,97.5)
\end{lstlisting}

\begin{bmcavbox}{Overall PS and the CI are computed two different ways}
The overall PS averages the components first, then weights them. The CI averages per-row composites. These are not the same operation --- for non-linear cases and for models with many empty/NaN rows they can diverge, and the overall PS can even sit outside its own CI (in the bundled run kimi's PS\_overall 0.435 lies above its CI [0.184, 0.247]). Read the CI as a stability indicator for the per-instance score, not as a strict interval around PS\_overall. If you need them to coincide, compute the point estimate the same per-row way.
\end{bmcavbox}

\subsubsection{The capability radar}

The \code{radar\_chart} plots the five [0,1] axes \textbf{FASR \(\cdot\) IWSR \(\cdot\) Exec \(\cdot\) IHR \(\cdot\) PAS} as a pentagon per model. It is the visual companion to the profile thresholds: a Genuine Planner fills most of the pentagon; a No-Grounding model collapses to the centre; a Vocabulary-Only model shows a lopsided shape --- high IHR but low PAS and Exec. Read \emph{area} for overall strength and \emph{shape} for the capability signature. Note the radar uses IHR (not 1-halluc) and omits the CoT bonus, so it is the \emph{base} profile, not the bonus-adjusted PS.

\begin{bmcavbox}{Radar reading caveat}
Equal area can come from very different shapes --- a model strong on FASR/IWSR but weak on Exec/PAS can match the area of a balanced model. Always compare shape, not just fill. With many models the overlapping fills become unreadable; facet or pick a handful for any figure you publish.
\end{bmcavbox}

\subsubsection{Goal-dependent weight tuning}

The weights encode a scientific claim; changing them changes the question PS answers. They live in one dict (\code{COMPOSITE\_WEIGHTS}) and must sum to 1.0. Reasonable retunings for different research goals:

\begin{table}[h]\small
\centering
\begin{tabularx}{\linewidth}{@{}l Y@{}}
\toprule
Research goal & Adjustment\\
\midrule
Zero-shot capability only & FASR\(\rightarrow\)0.40, IWSR\(\rightarrow\)0.10\\
\addlinespace
Domain grounding focus & one\_minus\_halluc\(\rightarrow\)0.35, FASR\(\rightarrow\)0.15\\
\addlinespace
All dimensions equal & all five = 0.20 (simple average)\\
\addlinespace
Execution-centric & exec\_ratio\(\rightarrow\)0.35, PAS\(\rightarrow\)0.10\\
\bottomrule
\end{tabularx}
\end{table}

\begin{bmkeybox}{The discipline that makes PS safe}
Always report the component scores next to PS. A PS of 0.60 from FASR=1.0 / exec=0.0 (always right on the first try, never produces a partially-executable failure) is a completely different model from a PS of 0.60 from moderate scores everywhere. The composite ranks; the components explain.
\end{bmkeybox}
